Docking systems are usually presented as a tradeoff: coarse approximations make docking faster, while richer scoring and more careful search make it more accurate. This paper studies a different regime in which the right kind of compression improves both. The key idea is decision-relevant locality. If approximation error stays below half the strict utility gap, then the optimizer is unchanged; if interactions outside a cutoff cannot perturb utility beyond that margin, then those coordinates are decision-irrelevant; and if only a small binding pocket can influence the optimizer, then the structural rank of the docking problem collapses from the ambient molecular dimension to a much smaller decision-relevant core. From this viewpoint, fast docking is not obtained by throwing away signal, but by discarding only information that has already been proved irrelevant to the docking decision. We formalize this theorem spine using Lean-backed approximation, locality, structural-rank, sampled-docking, and pruning results, and we implement it in DQDock, a theorem-guided virtual docking engine with explicit certified, conditionally certified, empirical, and heuristic layers. The empirical payoff is unusually strong: on a current suite of ten nontrivial protein--ligand complexes, including HIV-1 protease, ABL1, BACE-1, and SARS-CoV-2 main protease, DQDock averages below 1~\AA{} RMSD while a Vina baseline is around 5~\AA{}, and DQDock is roughly $20\times$ faster on average. The paper's claim is therefore not merely that formal structure can coexist with practical docking, but that decision-relevant compression can be the source of practical advantage.
